\documentclass[10pt,letterpaper]{article}

\usepackage[margin=0.58in,top=0.45in,bottom=0.45in]{geometry}
\usepackage[T1]{fontenc}
\usepackage{lmodern}
\usepackage{enumitem}
\usepackage{hyperref}
\usepackage{microtype}
\usepackage{titlesec}
\usepackage{parskip}

\hypersetup{colorlinks=true,urlcolor=black,linkcolor=black}
\setlength{\parindent}{0pt}
\setlength{\parskip}{0pt}
\setlist[itemize]{leftmargin=1.35em,itemsep=0.7pt,topsep=1pt,parsep=0pt,partopsep=0pt}

\titleformat{\section}{\bfseries\large}{}{0em}{}[\vspace{-2pt}\titlerule]
\titlespacing{\section}{0pt}{5pt}{2pt}

\newcommand{\role}[4]{%
  \textbf{#1} \hfill \textit{#2}\\
  \textit{#3} \hfill #4\vspace{-2pt}
}

\begin{document}

\begin{center}
{\LARGE \textbf{Yaze (Arthur) Li, Ph.D.}}\\[2pt]
{\small Power System Engineering Manager \textbar{} Interconnection Studies \textbar{} PSS\textregistered{}E \textbar{} PJM/MISO}\\[3pt]
{\small Calgary, AB, Canada \textbar{} +1 (672) 558-8667 \textbar{}
\href{mailto:arthurliyaze@gmail.com}{arthurliyaze@gmail.com}}\\[-1pt]
{\small \href{https://www.linkedin.com/in/yaze-li-6b436a228}{LinkedIn} \textbar{}
\href{https://github.com/Arthurliyaze}{GitHub} \textbar{}
\href{https://arthurliyaze.github.io}{Portfolio}}
\end{center}

\section{Professional Summary}
Power systems engineering manager delivering utility-scale interconnection studies across PJM, MISO, ERCOT, NYISO, and AESO. Expertise in power flow, short-circuit, transient stability, dynamic model validation, MQT, and automation using PSS\textregistered{}E, TARA, Aspen OneLiner, and Python. Experienced leading teams and converting engineering workflows into scalable tools that improve study speed and quality.

\section{Technical Skills}
\textbf{Power System Analysis:} PSS\textregistered{}E, TARA, Aspen OneLiner, TSAT/PSAT; power flow, short circuit, transient stability, voltage stability, VRT, MQT\\
\textbf{Interconnection Studies:} PJM, MISO, ERCOT, NYISO, AESO; Attachment M-3 DNH, affected system, replacement and surplus studies, cluster studies, case development, contingency definition, dynamic model validation\\
\textbf{PJM Platforms \& Standards:} Queue Point, NextGen, Dynamic Modeling Design Guide (DMDG)\\
\textbf{Programming \& Automation:} Python, MATLAB, SQL, R; PSS\textregistered{}E API, TARAViewer API, Pandas, NumPy, Git

\section{Professional Experience}

\role{Power System Engineering Manager}{Aug 2025 -- Present}{E Source}{Calgary, AB}
\begin{itemize}
  \item Lead PJM Attachment M-3 Do-No-Harm studies and develop Python automation using the TARAViewer API to evaluate Transmission Owner local transmission projects across AEP, Dominion, PPL, FirstEnergy, BGE, and PECO territories.
  \item Perform multiple rounds of PJM data review and validation through Queue Point and NextGen, verifying load-flow, short-circuit, and dynamic model data in accordance with the PJM Dynamic Modeling Design Guide (DMDG).
  \item Lead MISO Affected System Studies and Replacement/Surplus Interconnection Studies for utility-scale wind, solar, and battery energy storage projects.
  \item Perform PSS\textregistered{}E Model Quality Testing and dynamic model validation for NYISO projects in accordance with EMT Modeling Guideline UG28.
  \item Manage the steady-state study team and develop and maintain Python automation tools that streamline ERCOT interconnection study workflows for wind, solar, and BESS projects.
\end{itemize}

\role{Power System Engineer}{Aug 2024 -- Aug 2025}{CF Power Ltd. -- an E Source Company}{Calgary, AB}
\begin{itemize}
  \item Led steady-state interconnection studies for utility-scale renewable projects in ERCOT, including power flow, short-circuit, transient stability, and voltage ride-through assessments using PSS\textregistered{}E, TARA, and Aspen OneLiner.
  \item Supported PJM RTEP 2029 case development and performed dynamic simulation studies for TC1 and TC2 cluster projects.
  \item Executed PJM Cluster Studies including case development, generator model validation, contingency definition, dispatch preparation, and transient stability analysis.
  \item Developed Python-based PSS\textregistered{}E automation tools that reduced manual effort and improved consistency for AESO transient stability studies.
\end{itemize}

\role{Graduate Assistant}{Aug 2017 -- Jul 2023}{University of Arkansas (SEEDS \& CITES)}{Fayetteville, AR}
\begin{itemize}
  \item Developed dynamic programming and deep reinforcement learning scheduling algorithms for battery-assisted photovoltaic systems, reducing modeled costs by 29.3\% over a ten-year horizon.
  \item Built renewable-energy scheduling and cyberattack-detection frameworks using Pandapower, Stable Baselines, DQN, and dynamic watermarking.
\end{itemize}

\section{Education}
\textbf{University of Arkansas} \hfill \textit{August 2023}\\
Ph.D., Electrical Engineering \hfill Fayetteville, AR\\[2pt]
\textbf{Tsinghua University} \hfill \textit{July 2017}\\
B.S., Electrical, Electronics and Communication Engineering \hfill Beijing, China

\section{Selected Publications}
Seven peer-reviewed publications in venues including \textit{IEEE Transactions on Engineering Management}, \textit{International Journal of Electrical Power \& Energy Systems}, and \textit{IEEE JESTIE}. \href{https://scholar.google.com/citations?view_op=list_works&hl=en&user=WdrLhUkAAAAJ&sortby=pubdate}{Google Scholar}.

\end{document}
